% !TeX root = ../教材题库_全部条目_转写工作稿.tex
% Source: XJ-B1-C04-RV-P153-Y0875
% 第4章 / 章末复习 复习题四 / 教材 p.153 / PDF 第160页
% OCR未校对：公式与文字均可能有误；未完成数学TeX转写；禁止直接发给学生
% 原题及相邻原解截图（核对用）：../题目原图/XJ-B1-C04-RV-P153-Y0875-01.webp
% 本文件只在转写完成后作为题干使用；OCR 条目尚不是可靠题干。
\begin{bankitem}{复习题四 第22题}{第4章 / 章末复习 复习题四 / 教材 p.153 / PDF 第160页}{XJ-B1-C04-RV-P153-Y0875}{OCR未校对：公式与文字均可能有误；未完成数学TeX转写；禁止直接发给学生}
\begin{ocrdraft}
[教材p.153 / PDF第160页]
\par
22．
\par
（多选题）将一2的图象（），可得到函数3，一10g20十I)的图象．
\par
（A)先向左平移1个单位长度，再作关于直线3-工对称的图象
\par
(B)
\par
先作关于直线对称的图象，再向左平移1个单位长度
\par
（0先作关于直线3，一对称的图象，再向上平移1个单位长度
\par
(D)
\par
先向下平移1个单位长度，再作关于直线3-对称的图象
\end{ocrdraft}
\end{bankitem}
